\documentclass[10pt,a4paper]{amsart}
\usepackage[margin=19mm]{geometry}
\usepackage{amsmath,amssymb,mathtools}
\usepackage[scr=rsfs,cal=euler]{mathalfa}
\usepackage{xfrac}
\usepackage[unicode,hidelinks,bookmarks=false]{hyperref}
\newcommand{\ie}{i.e.\ }
\newcommand{\cf}{cf.\ }
\newcommand{\resp}{respectively\ }
\newcommand{\Subsection}{\subsection}
\providecommand{\nobreakdash}{\nobreak\hbox{-}}
\setlength{\emergencystretch}{2em}
\multlinegap0pt
\usepackage{mathtools}
\usepackage{enumitem}
\hypersetup{
  colorlinks=true,
  linkcolor=blue,
  citecolor=blue,
  urlcolor=blue
}
\newtheorem{theorem}{Theorem}[section]
\newtheorem{lemma}[theorem]{Lemma}
\newtheorem{corollary}[theorem]{Corollary}
\newtheorem{proposition}[theorem]{Proposition}
\newtheorem{remark}[theorem]{Remark}
\newcommand{\C}{\mathbb{C}}
\newcommand{\R}{\mathbb{R}}
\newcommand{\E}{\mathbb{E}}
\newcommand{\Et}{\widetilde{\mathbb{E}}}
\newcommand{\ip}[2]{\left\langle #1,#2\right\rangle}
\newcommand{\Tr}{\operatorname{Tr}}
\newcommand{\Herm}{\operatorname{Herm}}
\newcommand{\diag}{\operatorname{diag}}
\newcommand{\rank}{\operatorname{rank}}
\begin{document}
\title[Degree-Four Vector-Coordinate SoS Cannot Detect the MUB Uppe]{Degree-Four Vector-Coordinate SoS Cannot Detect the MUB Upper Bound}
\author{Shreyhaan Sarkar}
\date{}
\maketitle
\noindent\emph{Source and licence.} Degree-Four Vector-Coordinate SoS Cannot Detect the MUB Upper Bound. Version arXiv:2606.13903v1 (CC BY 4.0). This material is distributed under the Creative Commons Attribution 4.0 International licence, http://creativecommons.org/licenses/by/4.0/, as recorded in the arXiv OAI licence metadata for this exact version and shown on the abstract page https://arxiv.org/abs/2606.13903v1. Full rights notice: licenses/arxiv-2606.13903-CC-BY-4.0.txt.\par\smallskip
\noindent\textit{Editorial scope.} Counted unit: the original \texttt{theorem} environment \texttt{thm:affine-multiplier} at source lines 281--294 together with its complete proof at lines 296--421; the delivered document keeps source lines 273--422 so that every referenced label and every citation resolves inside it. Native editable source is the paper's own concatenated TeX; the AMS wrapper is editorial formatting only. No publisher class, upstream script or unrelated artwork is redistributed.
\section{An affine Haar multiplier theorem}

We now isolate the estimate needed for the $2\times2$ Hermitian LMI formulation. Let $x,y$ be independent Haar-random unit vectors in $\C^d$, and set
\[
z=\ip{x}{y}.
\]
The theorem says that multiplication by $z$ couples affine-linear functions only weakly: its operator norm is at most $1/d$.


\begin{theorem}[Affine Haar multiplier bound]\label{thm:affine-multiplier}
Let $x,y$ be independent Haar-random unit vectors in $\C^d$, and let $z=\ip{x}{y}$. Let $\mathcal A$ be the complex vector space of affine-linear polynomials in the vector-coordinate variables of any finite collection of independent Haar-random orthonormal bases containing $x$ and $y$. Then for all $p,q\in\mathcal A$,
\[
\left|\E[\overline p\,zq]\right|
\le
\frac1d
\sqrt{\E|p|^2\,\E|q|^2}.
\]
Equivalently, the bilinear form
\[
(p,q)\longmapsto \E[\overline p\,\ip{x}{y}\,q]
\]
has operator norm at most $1/d$ with respect to the $L^2$ norms on affine-linear functions.
\end{theorem}

\begin{proof}
We use the convention
\[
\ip{x}{y}=\sum_{r=1}^d \overline{x_r}y_r.
\]
The joint law of the Haar bases is invariant under independent phase rotations of each column. In particular, it is invariant under
\[
x\mapsto e^{i\theta}x,
\qquad
 y\mapsto e^{i\phi}y,
\]
and under independent phase rotations of all other columns.

Complexify the affine-linear space. Any affine-linear polynomial is a sum of a constant, linear coordinate terms, and conjugate-linear coordinate terms. In
\[
\E[\overline p\,zq],
\]
a monomial contribution can be nonzero only if its total phase is neutral under every independent column phase.

The phase bookkeeping is as follows; each row denotes any coordinate of the indicated vector:
\[
\begin{array}{c|c}
\text{factor} & \text{phase under } x\mapsto e^{i\theta}x,\; y\mapsto e^{i\phi}y \\
\hline
x & e^{i\theta}\\
\overline{x} & e^{-i\theta}\\
y & e^{i\phi}\\
\overline{y} & e^{-i\phi}\\
z=\ip{x}{y} & e^{-i\theta}e^{i\phi}
\end{array}
\]
A monomial contribution in $\E[\overline p\,zq]$ survives only if its total phase is trivial under both independent phase rotations. If a term of $p$ has phase exponent $\alpha\in\mathbb Z^2$ and a term of $q$ has phase exponent $\beta\in\mathbb Z^2$, then the corresponding term in $\overline p\,zq$ has exponent
\[
-\alpha+(-1,1)+\beta.
\]
Thus phase balance requires
\[
\beta=\alpha+(1,-1).
\]
Because $p$ and $q$ are affine-linear, the only possible matches are
\[
\alpha=(0,1),\quad \beta=(1,0),
\]
which gives
\[
p=a\cdot y,
\qquad
q=b\cdot x,
\]
and
\[
\alpha=(-1,0),\quad \beta=(0,-1),
\]
which gives
\[
p=a\cdot \overline x,
\qquad
q=b\cdot \overline y,
\]
for vectors $a,b\in\C^d$. Constant terms, all terms involving other columns, and all other $x,y$-linear or conjugate-linear combinations have unbalanced phase and integrate to zero.

For the first block, write
\[
p=a\cdot y=\sum_{r=1}^d a_r y_r,
\qquad
q=b\cdot x=\sum_{t=1}^d b_t x_t.
\]
Then
\[
\overline p=\sum_{r=1}^d \overline{a_r}\,\overline{y_r},
\qquad
z=\sum_{s=1}^d \overline{x_s}y_s,
\qquad
q=\sum_{t=1}^d b_t x_t.
\]
Using independence of $x$ and $y$,
\[
\begin{aligned}
\E[\overline p\,zq]
&=
\sum_{r,s,t}
\overline{a_r}b_t
\E[\overline{y_r}y_s]
\E[\overline{x_s}x_t].
\end{aligned}
\]
For a Haar-random unit vector $u\in\C^d$,
\[
\E[\overline{u_r}u_s]=\frac{\delta_{rs}}d.
\]
Therefore
\[
\E[\overline p\,zq]
=
\frac1{d^2}\sum_{r=1}^d \overline{a_r}b_r.
\]
Also,
\[
\E|p|^2=\frac{\|a\|^2}{d},
\qquad
\E|q|^2=\frac{\|b\|^2}{d}.
\]
Thus
\[
\left|\E[\overline p\,zq]\right|
\le
\frac{\|a\|\|b\|}{d^2}
=
\frac1d\sqrt{\E|p|^2\,\E|q|^2}.
\]

The second nonzero block, $p=a\cdot\overline x$ and $q=b\cdot\overline y$, gives the same estimate. The two blocks are orthogonal in $L^2$ by phase invariance. More explicitly, if $p_y$ and $p_{\overline x}$ denote the projections of $p$ onto the spans of the coordinates of $y$ and of $\overline x$, respectively, and $q_x,q_{\overline y}$ are defined similarly, then
\[
\E[\overline p\,zq]
=
\frac1{d^2}
\left(\langle a_y,b_x\rangle+\langle a_{\overline x},b_{\overline y}\rangle\right),
\]
whereas
\[
\E|p|^2\ge \frac{\|a_y\|^2+\|a_{\overline x}\|^2}{d},
\qquad
\E|q|^2\ge \frac{\|b_x\|^2+\|b_{\overline y}\|^2}{d}.
\]
Cauchy's inequality on the direct sum of the two relevant blocks gives the claimed bound.
\end{proof}


\end{document}
